% ---------------------------------------------------------------------
% Bibliografia
%
% Composta a mano e per sezioni, come vuole l'uso giuridico italiano.
% NON si usa BibTeX: il sistema autore-anno e gli stili bibliografici
% standard (plain, APA, IEEE) sono estranei a questo tipo di lavoro.
%
% Regole:
%  - ogni voce qui elencata deve avere una riga in output/registro-fonti.md;
%  - ogni voce deve essere citata almeno una volta in nota;
%  - ordine alfabetico per autore nella dottrina, cronologico nelle altre
%    sezioni;
%  - lo stile delle voci segue
%    .claude/skills/tesi-contabilita-pubblica/references/metodo-citazione.md.
%
% Controllo automatico: python3 scripts/audit_note_tesi.py
% ---------------------------------------------------------------------

\chapter*{Bibliografia}
\addcontentsline{toc}{chapter}{Bibliografia}

\section*{Fonti normative}
\addcontentsline{toc}{section}{Fonti normative}
\begin{itemize}\setlength{\itemsep}{0.2em}
  \item % legge 7 gennaio 2026, n. 1, ...
\end{itemize}

\section*{Giurisprudenza costituzionale e di legittimità}
\addcontentsline{toc}{section}{Giurisprudenza costituzionale e di legittimità}
\begin{itemize}\setlength{\itemsep}{0.2em}
  \item % Corte cost., sent. ..., n. ...
\end{itemize}

\section*{Deliberazioni e pronunce della Corte dei conti}
\addcontentsline{toc}{section}{Deliberazioni e pronunce della Corte dei conti}
\begin{itemize}\setlength{\itemsep}{0.2em}
  \item % Corte conti, sez. riun. contr., delib. n. .../QMIG
\end{itemize}

\section*{Giurisprudenza amministrativa}
\addcontentsline{toc}{section}{Giurisprudenza amministrativa}
\begin{itemize}\setlength{\itemsep}{0.2em}
  \item % Cons. St., sez. ..., sent. ..., n. ...
\end{itemize}

\section*{Dottrina}
\addcontentsline{toc}{section}{Dottrina}
\begin{itemize}\setlength{\itemsep}{0.2em}
  \item % A. Rossi, \emph{Titolo}, Milano, 2024.
\end{itemize}

\section*{Atti parlamentari e documentazione}
\addcontentsline{toc}{section}{Atti parlamentari e documentazione}
\begin{itemize}\setlength{\itemsep}{0.2em}
  \item % Camera dei deputati, Servizio studi, dossier n. ..., data.
\end{itemize}
